\makeatletter
\def\input@path{{./}{../}{../../}{../../../}{../../../../}{preamble/}{../preamble/}{../../preamble/}{../../../preamble/}{../../../../preamble/}{book2026-2027/preamble/}{book2026-2027/}}
\makeatother

\input{preamble_G10_T1_C1C2C3}
\input{preamble_exam_layout}

\title{Decision-Making Under Uncertainty: Exam Solutions}
\author{}
\date{}

\begin{document}
	\maketitle
	\vspace{-2cm}
	\setcounter{tocdepth}{1}
	\setcounter{secnumdepth}{2}
	\phantomsection
	\label{toc}
\label{sec:solutions-extraction}
	\tableofcontents
	
	\vspace{-5mm}
	\section*{Learning objective}\vspace{-5mm}
	Structure decision-making problems under uncertainty, calculate and organise payoffs, and apply the Maximax and Maximin criteria to compare alternatives and justify decisions.
	
	\section*{Criteria assessment}\vspace{-5mm}
	This task sheet assesses criteria C1 through C4. Each problem assesses the criteria indicated in its title. The number of points assigned to each task is shown at the end of the item. Partial credit may be awarded when the correct procedure is shown but arithmetic errors are made. Answers that do not show the required procedure will not receive credit. A minimum of 4 points is required to pass each criterion.
	
	\section*{Evaluated criteria}
	\begin{enumerate}[label=\textbf{C\arabic*}]
		\item Characterises decision-making problems under uncertainty by identifying decision alternatives, states of nature, consequences, and the economic or financial context of the decision.
		\item Organises decision problems using payoff tables, placing alternatives, states of nature, and payoffs in a clear structure for comparison.
		\item Calculates profits by combining revenues with fixed and variable costs across different states of nature, determining the corresponding payoffs for each alternative, and organising the results in payoff tables for comparison.
		\item Applies the Maximin and Maximax criteria to make optimistic or conservative decisions without probabilities, selecting the alternative with the highest possible payoff or the best worst-case, interpreting risks or limitations.
	\end{enumerate}
	
	\sectionScored
	{Local Food Market Stall}
	{C1-\sml{1} \hspace{2mm} C2-\sml{2}}
	{sec:local-food-market-stall-c1-c2}
	{
		A local food market company is planning its stall offering for an upcoming market period. A standard stall is expected to generate profits of eighteen thousand US dollars under high demand and eleven thousand US dollars under low demand, while a deluxe stall is expected to generate thirty-one thousand US dollars and six thousand US dollars, respectively. The company must select its stall before the actual level of customer demand is known.
		
		\textbf{Questions/Tasks.}
		\begin{enumerate}
			\CriterionTask
			{subsec:local-food-market-stall-c1-c2-c1-characterization}{1}
			{Characterize the key elements of the decision-making situation. \score{1}}
			
			\CriterionTask
			{subsec:local-food-market-stall-c1-c2-c2-payoff-table}{2}
			{Construct the payoff table. \score{2}}
		\end{enumerate}
		
		\phantomsection
		\label{subsec:local-food-market-stall-c1-c2-c1-characterization}
		\subsection*{C1 - Characterization \score{1}}
		
		\DecisionElementsTable{Standard stall; Deluxe stall}{High demand; Low demand}{Actual market demand after the stall decision}{Operating the selected activity under the realised conditions leaves funds available or a shortfall to cover. Monetary measure: Profit in thousand US dollars for each alternative--state pair}
		
		\phantomsection
		\label{subsec:local-food-market-stall-c1-c2-c2-payoff-table}
		\subsection*{C2 - Payoff table \score{2}}
		
		Each supplied monetary payoff is placed at the intersection of its alternative and state:
		\[
		\begin{array}{ll}
			P(A_1,S_1)=18, & P(A_1,S_2)=11,\\
			P(A_2,S_1)=31, & P(A_2,S_2)=6.
		\end{array}
		\]
		Thus, in thousand US dollars, the completed payoff table is
		\begin{center}
			\begin{tabular}{lcc}
				\toprule
				Alternative & High \(S_1\) & Low \(S_2\)\\
				\midrule
				Standard stall \(A_1\) & 18 & 11\\
				Deluxe stall \(A_2\) & 31 & 6\\
				\bottomrule
			\end{tabular}
		\end{center}
		
		\vspace{-5mm}
		\secInicio[sec:local-food-market-stall-c1-c2][subsec:local-food-market-stall-c1-c2-c2-payoff-table]
		
	}
	
	\sectionScored
	{Small Bakery Production}
	{C1-\sml{1} \hspace{2mm} C2-\sml{3}}
	{sec:small-bakery-production-c1-c2}
	{
		A small bakery is planning its production capacity for an upcoming sales period. A standard oven can generate 3 thousand BRL per production batch, while an industrial oven can generate 5 thousand BRL per batch. Based on expected demand, the standard oven would complete 16 batches under high demand or 11 under low demand, whereas the industrial oven would complete 13 or 9 batches, respectively. In addition, the bakery would receive a fixed sales bonus of 6 thousand BRL if demand is high and incur a fixed loss of 3 thousand BRL if demand is low.
		
		\textbf{Questions/Tasks.}
		\begin{enumerate}
			\CriterionTask
			{subsec:small-bakery-production-c1-c2-c1-characterization}{1}
			{Characterize the key elements of the decision-making situation. \score{1}}
			
			\CriterionTask
			{subsec:small-bakery-production-c1-c2-c2-payoff-table}{2}
			{Construct the payoff table. \score{3}}
		\end{enumerate}
		
		\phantomsection
		\label{subsec:small-bakery-production-c1-c2-c1-characterization}
		\subsection*{C1 - Characterization \score{1}}
		
		\DecisionElementsTable
{Standard oven; Industrial oven}
{High demand; Low demand}
{Actual customer demand after the oven decision}
{Operating the selected activity under the realised conditions leaves funds available or a shortfall to cover. The monetary measure is profit in thousand BRL for each alternative--state pair.}

		
		\begin{CalculationLandscape}{11in}{8.5in}
\phantomsection
		\label{subsec:small-bakery-production-c1-c2-c2-payoff-table}
		\subsection*{C2 - Payoff table \score{3}}
		
		
\SymbolsTable{
$A_i$ & $A_1$: Standard oven; $A_2$: Industrial oven & Alternative labels, in row order. \\
$S_j$ & $S_1$: High demand; $S_2$: Low demand & State labels, in column order; revealed after the choice. \\
$q(A_i,S_j)$ & Standard--high: 16 batches; Standard--low: 11 batches; Industrial--high: 13 batches; Industrial--low: 9 batches & Expected completed production batches for each oven under each demand condition. \\
$v^A(A_i)$ & Standard oven: 3 thousand BRL/batch; Industrial oven: 5 thousand BRL/batch & Earnings per completed production batch from the chosen oven. \\
$f^S(S_j)$ & High demand: 6 thousand BRL; Low demand: $-3$ thousand BRL & Fixed sales bonus or loss associated with the demand condition. \\
}

		
		\textbf{Payoff equation.}
		\[
		P_{ay}(A_i,S_j)
		=
		q(A_i,S_j)v^A(A_i)
		+
		f^S(S_j).
		\]
		
		\CalculationPair{Additive payoff}{
\begin{tabular}{@{}lcc@{}}
\toprule Alternative & $S_1$ & $S_2$ \\ \midrule
$A_1$ & \SymbolicCell{q(A_1,S_1)\times v^A(A_1) \\ {}+f^S(S_1)}{P_{ay}(A_1,S_1)} & \SymbolicCell{q(A_1,S_2)\times v^A(A_1) \\ {}+f^S(S_2)}{P_{ay}(A_1,S_2)} \\[5pt]
$A_2$ & \SymbolicCell{q(A_2,S_1)\times v^A(A_2) \\ {}+f^S(S_1)}{P_{ay}(A_2,S_1)} & \SymbolicCell{q(A_2,S_2)\times v^A(A_2) \\ {}+f^S(S_2)}{P_{ay}(A_2,S_2)} \\[5pt]
\bottomrule
\end{tabular}
}{
\begin{tabular}{@{}lcc@{}}
\toprule Alternative & $S_1$ & $S_2$ \\ \midrule
$A_1$ & \CalcResult{16\times 3 \\ {}+6}{54} & \CalcResult{11\times 3 \\ {}-3}{30} \\[5pt]
$A_2$ & \CalcResult{13\times 5 \\ {}+6}{71} & \CalcResult{9\times 5 \\ {}-3}{42} \\[5pt]
\bottomrule
\end{tabular}
}
Negative rates reduce the per-unit contribution; negative fixed amounts represent losses or commitments deducted once. A negative final payoff is a shortfall.

		
		\begingroup\renewcommand{\CalculationStateCount}{2}
\FinalPayoffTable{profit in thousand BRL}{High demand ($S_1$) & Low demand ($S_2$)}{
Standard oven ($A_1$) & 54 & 30 \\
Industrial oven ($A_2$) & 71 & 42 \\
}
\endgroup
\end{CalculationLandscape}

		
		\vspace{-5mm}
		\secInicio[sec:small-bakery-production-c1-c2][subsec:small-bakery-production-c1-c2-c2-payoff-table]
	}
	
	\sectionScored
	{Urban Meal Subscription Service}
	{C1-\sml{2} \hspace{2mm} C2-\sml{5}}
	{sec:urban-meal-subscription-service-c1-c2}
	{
		An urban meal subscription company must choose between a basic meal plan and a premium meal plan before customer interest is known. Interest may be strong or weak. The basic plan generates 4 thousand BRL per active subscriber and also provides a fixed operating contribution of 6 thousand BRL. The premium plan generates 6 thousand BRL per active subscriber and provides a fixed operating contribution of 9 thousand BRL. Strong interest adds another 2 thousand BRL per active subscriber to either plan, while weak interest adds 1 thousand BRL per active subscriber. The basic plan is expected to have 13 active subscribers under strong interest and 8 under weak interest; the premium plan is expected to have 10 active subscribers under strong interest and 6 under weak interest.
		
		\textbf{Questions/Tasks.}
		\begin{enumerate}
			\CriterionTask
			{subsec:urban-meal-subscription-service-c1-c2-c1-characterization}{1}
			{Characterize the key elements of the decision-making situation. \score{2}}
			
			\CriterionTask
			{subsec:urban-meal-subscription-service-c1-c2-c2-payoff-table}{2}
			{Construct the payoff table. \score{5}}
		\end{enumerate}
		
		\phantomsection
		\label{subsec:urban-meal-subscription-service-c1-c2-c1-characterization}
		\subsection*{C1 - Characterization \score{2}}
		
		\DecisionElementsTable
{Basic meal plan; Premium meal plan}
{Strong customer interest; Weak customer interest}
{Actual customer interest after the meal plan is selected}
{Operating the selected activity under the realised conditions leaves funds available or a shortfall to cover. The monetary measure is profit in thousand BRL for each alternative--state pair.}

		
		\begin{CalculationLandscape}{11in}{8.5in}
\phantomsection
		\label{subsec:urban-meal-subscription-service-c1-c2-c2-payoff-table}
		\subsection*{C2 - Payoff table \score{5}}
		
		
\SymbolsTable{
$A_i$ & $A_1$: Basic meal plan; $A_2$: Premium meal plan & Alternative labels, in row order. \\
$S_j$ & $S_1$: Strong customer interest; $S_2$: Weak customer interest & State labels, in column order; revealed after the choice. \\
$q(A_i,S_j)$ & Basic--strong: 13 subscribers; Basic--weak: 8 subscribers; Premium--strong: 10 subscribers; Premium--weak: 6 subscribers & Expected active subscribers for each meal plan under each level of customer interest. \\
$v^A(A_i)$ & Basic plan: 4 thousand BRL/subscriber; Premium plan: 6 thousand BRL/subscriber & Contribution per active subscriber from the chosen meal plan. \\
$v^S(S_j)$ & Strong interest: 2 thousand BRL/subscriber; Weak interest: 1 thousand BRL/subscriber & Additional contribution per active subscriber from customer interest. \\
$f^A(A_i)$ & Basic plan: 6 thousand BRL; Premium plan: 9 thousand BRL & Fixed operating contribution from the chosen meal plan. \\
}

		
		\textbf{Payoff equation.}
		\[
		P_{ay}(A_i,S_j)
		=
		q(A_i,S_j)
		\left[
		v^A(A_i)+v^S(S_j)
		\right]
		+
		f^A(A_i).
		\]
		
		\CalculationPair{Additive payoff}{
\begin{tabular}{@{}lcc@{}}
\toprule Alternative & $S_1$ & $S_2$ \\ \midrule
$A_1$ & \SymbolicCell{q(A_1,S_1)\times \\ [v^A(A_1)+v^S(S_1)] \\ {}+f^A(A_1)}{P_{ay}(A_1,S_1)} & \SymbolicCell{q(A_1,S_2)\times \\ [v^A(A_1)+v^S(S_2)] \\ {}+f^A(A_1)}{P_{ay}(A_1,S_2)} \\[5pt]
$A_2$ & \SymbolicCell{q(A_2,S_1)\times \\ [v^A(A_2)+v^S(S_1)] \\ {}+f^A(A_2)}{P_{ay}(A_2,S_1)} & \SymbolicCell{q(A_2,S_2)\times \\ [v^A(A_2)+v^S(S_2)] \\ {}+f^A(A_2)}{P_{ay}(A_2,S_2)} \\[5pt]
\bottomrule
\end{tabular}
}{
\begin{tabular}{@{}lcc@{}}
\toprule Alternative & $S_1$ & $S_2$ \\ \midrule
$A_1$ & \CalcResult{13\times \\ [4+2] \\ {}+6}{84} & \CalcResult{8\times \\ [4+1] \\ {}+6}{46} \\[5pt]
$A_2$ & \CalcResult{10\times \\ [6+2] \\ {}+9}{89} & \CalcResult{6\times \\ [6+1] \\ {}+9}{51} \\[5pt]
\bottomrule
\end{tabular}
}

		
		\begingroup\renewcommand{\CalculationStateCount}{2}
\FinalPayoffTable{profit in thousand BRL}{Strong customer interest ($S_1$) & Weak customer interest ($S_2$)}{
Basic plan ($A_1$) & 84 & 46 \\
Premium plan ($A_2$) & 89 & 51 \\
}
\endgroup
\end{CalculationLandscape}

		
		\vspace{-5mm}
		\secInicio[sec:urban-meal-subscription-service-c1-c2][subsec:urban-meal-subscription-service-c1-c2-c2-payoff-table]
		
	}
	
	\sectionScored
	{Regional Bookstore Expansion}
	{C1-\sml{1} \hspace{2mm} C3-\sml{2}}
	{sec:regional-bookstore-expansion-c1-c3}
	{
		A regional bookstore must decide whether to open a children's section or maintain its current product selection before customer traffic is known. Traffic may be strong or weak. All values are in million BRL. Opening the children's section produces revenue of 72 under strong traffic and 35 under weak traffic, while maintaining the current selection produces revenue of 61 under strong traffic and 47 under weak traffic. The corresponding costs are 24 and 22 for the children's section under strong and weak traffic, respectively, and 25 and 23 for the current selection under strong and weak traffic, respectively.
		
		\textbf{Questions/Tasks.}
		\begin{enumerate}
			\CriterionTask
			{subsec:regional-bookstore-expansion-c1-c3-c1-characterization}{1}
			{Characterize the key elements of the decision-making situation. \score{1}}
			
			\CriterionTask
			{subsec:regional-bookstore-expansion-c1-c3-c3-profit-payoff-table}{3}
			{Use revenue and costs to construct the payoff table. \score{2}}
		\end{enumerate}
		
		\phantomsection
		\label{subsec:regional-bookstore-expansion-c1-c3-c1-characterization}
		\subsection*{C1 - Characterization \score{1}}
		
		\DecisionElementsTable
{Open children's section; Maintain current selection}
{Strong traffic; Weak traffic}
{Actual customer traffic after the product selection decision}
{Operating the selected activity under the realised conditions leaves funds available or a shortfall to cover. The monetary measure is profit in million BRL for each alternative--state pair.}

		
		\begin{CalculationLandscape}{11in}{8.5in}
\phantomsection
		\label{subsec:regional-bookstore-expansion-c1-c3-c3-profit-payoff-table}
		\subsection*{C3 - Calculating profit and building the payoff table \score{2}}
		
		
\SymbolsTable{
$A_i$ & $A_1$: Open children's section; $A_2$: Maintain current selection & Alternative labels, in row order. \\
$S_j$ & $S_1$: Strong traffic; $S_2$: Weak traffic & State labels, in column order; revealed after the choice. \\
$R(A_i,S_j)$ & Children's section--strong: 72 million BRL; Children's section--weak: 35 million BRL; Current selection--strong: 61 million BRL; Current selection--weak: 47 million BRL & Total revenue generated by each product selection under each traffic condition. \\
$C(A_i,S_j)$ & Children's section--strong: 24 million BRL; Children's section--weak: 22 million BRL; Current selection--strong: 25 million BRL; Current selection--weak: 23 million BRL & Total cost associated with each product selection under each traffic condition. \\
}

		
		\textbf{Profit equation.}
		\[
		P(A_i,S_j)
		=
		R(A_i,S_j)-C(A_i,S_j).
		\]
		
		\CalculationPair{Profit}{
\begin{tabular}{@{}lcc@{}}
\toprule Alternative & $S_1$ & $S_2$ \\ \midrule
$A_1$ & \SymbolicCell{R(A_1,S_1)-C(A_1,S_1)}{P(A_1,S_1)} & \SymbolicCell{R(A_1,S_2)-C(A_1,S_2)}{P(A_1,S_2)} \\[5pt]
$A_2$ & \SymbolicCell{R(A_2,S_1)-C(A_2,S_1)}{P(A_2,S_1)} & \SymbolicCell{R(A_2,S_2)-C(A_2,S_2)}{P(A_2,S_2)} \\[5pt]
\bottomrule
\end{tabular}
}{
\begin{tabular}{@{}lcc@{}}
\toprule Alternative & $S_1$ & $S_2$ \\ \midrule
$A_1$ & \CalcResult{72-24}{48} & \CalcResult{35-22}{13} \\[5pt]
$A_2$ & \CalcResult{61-25}{36} & \CalcResult{47-23}{24} \\[5pt]
\bottomrule
\end{tabular}
}

		
		\begingroup\renewcommand{\CalculationStateCount}{2}
\FinalPayoffTable{profit in million BRL}{Strong traffic ($S_1$) & Weak traffic ($S_2$)}{
Open children's section ($A_1$) & 48 & 13 \\
Maintain current selection ($A_2$) & 36 & 24 \\
}
\endgroup
\end{CalculationLandscape}

		
		\vspace{-5mm}
		\secInicio[sec:regional-bookstore-expansion-c1-c3][subsec:regional-bookstore-expansion-c1-c3-c3-profit-payoff-table]
		
	}
	
	
	\sectionScored
	{Greenhouse Crop Production}
	{C1-\sml{1} \hspace{2mm} C3-\sml{3}}
	{sec:greenhouse-crops-c1-c3}
	{
		A greenhouse farm is planning its production for an upcoming growing
		season. Revenue per production lot is 18 thousand EUR for tomatoes and
		15 thousand EUR for peppers. Under favourable weather conditions, the
		farm expects to produce 8 lots of tomatoes and 7 lots of peppers, while
		under unfavourable weather conditions it expects to produce 5 lots of
		tomatoes and 4 lots of peppers. The cost per production lot is 7 thousand
		EUR under favourable weather and 10 thousand EUR under unfavourable
		weather.
		
		\textbf{Questions/Tasks.}
		\begin{enumerate}
			\CriterionTask
			{subsec:greenhouse-crops-c1-c3-c1-characterization}{1}
			{Characterize the key elements of the decision-making situation. \score{1}}
			
			\CriterionTask
			{subsec:greenhouse-crops-c1-c3-c3-profit-table}{3}
			{Use revenue and costs to construct the payoff table. \score{3}}
		\end{enumerate}
		
		\phantomsection
		\label{subsec:greenhouse-crops-c1-c3-c1-characterization}
		\subsection*{C1 - Characterization \score{1}}
		
		\DecisionElementsTable
{Produce tomatoes; Produce peppers}
{Favourable weather; Unfavourable weather}
{Weather conditions revealed after the production decision}
{Operating the selected activity under the realised conditions leaves funds available or a shortfall to cover. The monetary measure is profit in thousand EUR for each crop--weather-state pair.}

		
		\begin{CalculationLandscape}{14in}{11in}
\phantomsection
		\label{subsec:greenhouse-crops-c1-c3-c3-profit-table}
		\subsection*{C3 - Profit and payoff table \score{3}}
		
		\secInicio[sec:solutions-extraction][sec:greenhouse-crops-c1-c3]
		

\SymbolsTable{
$A_i$ & $A_1$: Produce tomatoes; $A_2$: Produce peppers & Alternative labels, in row order. \\
$S_j$ & $S_1$: Favourable weather; $S_2$: Unfavourable weather & State labels, in column order; revealed after the choice. \\
$q(A_i,S_j)$ & Tomatoes--favourable: 8; Tomatoes--unfavourable: 5; Peppers--favourable: 7; Peppers--unfavourable: 4 & Production lots for each crop under each weather condition. \\
$r(A_i)$ & Tomatoes: 18 thousand EUR/lot; Peppers: 15 thousand EUR/lot & Revenue per production lot. \\
$c(S_j)$ & Favourable weather: 7 thousand EUR/lot; Unfavourable weather: 10 thousand EUR/lot & Cost per production lot. \\
}

		
		\textbf{Profit equation.}
		\[
		R=q(A_i,S_j)r(A_i),
		\qquad
		C=q(A_i,S_j)c(S_j),
		\qquad
		P(A_i,S_j)=R-C.
		\]
		
		\CalculationPair{Revenue}{
\begin{tabular}{@{}lcc@{}}
\toprule Alternative & $S_1$ & $S_2$ \\ \midrule
$A_1$ & \SymbolicCell{q(A_1,S_1)\times r(A_1)}{R(A_1,S_1)} & \SymbolicCell{q(A_1,S_2)\times r(A_1)}{R(A_1,S_2)} \\[5pt]
$A_2$ & \SymbolicCell{q(A_2,S_1)\times r(A_2)}{R(A_2,S_1)} & \SymbolicCell{q(A_2,S_2)\times r(A_2)}{R(A_2,S_2)} \\[5pt]
\bottomrule
\end{tabular}
}{
\begin{tabular}{@{}lcc@{}}
\toprule Alternative & $S_1$ & $S_2$ \\ \midrule
$A_1$ & \CalcResult{8\times 18}{144} & \CalcResult{5\times 18}{90} \\[5pt]
$A_2$ & \CalcResult{7\times 15}{105} & \CalcResult{4\times 15}{60} \\[5pt]
\bottomrule
\end{tabular}
}
\CalculationPair{Cost}{
\begin{tabular}{@{}lcc@{}}
\toprule Alternative & $S_1$ & $S_2$ \\ \midrule
$A_1$ & \SymbolicCell{q(A_1,S_1)\times c(S_1)}{C(A_1,S_1)} & \SymbolicCell{q(A_1,S_2)\times c(S_2)}{C(A_1,S_2)} \\[5pt]
$A_2$ & \SymbolicCell{q(A_2,S_1)\times c(S_1)}{C(A_2,S_1)} & \SymbolicCell{q(A_2,S_2)\times c(S_2)}{C(A_2,S_2)} \\[5pt]
\bottomrule
\end{tabular}
}{
\begin{tabular}{@{}lcc@{}}
\toprule Alternative & $S_1$ & $S_2$ \\ \midrule
$A_1$ & \CalcResult{8\times 7}{56} & \CalcResult{5\times 10}{50} \\[5pt]
$A_2$ & \CalcResult{7\times 7}{49} & \CalcResult{4\times 10}{40} \\[5pt]
\bottomrule
\end{tabular}
}
\CalculationPair{Profit}{
\begin{tabular}{@{}lcc@{}}
\toprule Alternative & $S_1$ & $S_2$ \\ \midrule
$A_1$ & \SymbolicCell{R(A_1,S_1)-C(A_1,S_1)}{P(A_1,S_1)} & \SymbolicCell{R(A_1,S_2)-C(A_1,S_2)}{P(A_1,S_2)} \\[5pt]
$A_2$ & \SymbolicCell{R(A_2,S_1)-C(A_2,S_1)}{P(A_2,S_1)} & \SymbolicCell{R(A_2,S_2)-C(A_2,S_2)}{P(A_2,S_2)} \\[5pt]
\bottomrule
\end{tabular}
}{
\begin{tabular}{@{}lcc@{}}
\toprule Alternative & $S_1$ & $S_2$ \\ \midrule
$A_1$ & \CalcResult{144-56}{88} & \CalcResult{90-50}{40} \\[5pt]
$A_2$ & \CalcResult{105-49}{56} & \CalcResult{60-40}{20} \\[5pt]
\bottomrule
\end{tabular}
}

		
		\begingroup\renewcommand{\CalculationStateCount}{2}
\FinalPayoffTable{profit in thousand EUR}{Favourable weather ($S_1$) & Unfavourable weather ($S_2$)}{
Produce tomatoes ($A_1$) & 88 & 40 \\
Produce peppers ($A_2$) & 56 & 20 \\
}
\endgroup
\end{CalculationLandscape}

		
		
	}

	\sectionScored
	{Fitness Center Membership Plans}
	{C1-\sml{2} \hspace{2mm} C3-\sml{5}}
	{sec:fitness-membership-c1-c3}
	{
		A fitness center must choose between a basic membership plan and a
		premium membership plan before monthly demand is known. Demand may be
		high or low. Quantities, in hundreds of memberships, are 9 and 6
		for the basic plan and 7 and 4 for the premium plan under high and
		low demand, respectively. The plan-based revenue rates are 8 and 13
		thousand USD per hundred memberships for the basic and premium
		plans. High demand adds a state-based revenue rate of 2 thousand USD
		per hundred memberships; low demand adds 1. The plan-based cost rates
		are 3 and 6 thousand USD per hundred memberships, respectively. The
		fixed operating costs are 12 thousand USD under high demand and
		16 thousand USD under low demand.
		
		\textbf{Questions/Tasks.}
		\begin{enumerate}
			\CriterionTask
			{subsec:fitness-membership-c1-c3-c1-characterization}{1}
			{Characterize the key elements of the decision-making situation. \score{2}}
			
			\CriterionTask
			{subsec:fitness-membership-c1-c3-c3-profit-table}{3}
			{Use revenue and costs to construct the payoff table. \score{5}}
		\end{enumerate}
		
		\phantomsection
		\label{subsec:fitness-membership-c1-c3-c1-characterization}
		\subsection*{C1 - Characterization \score{2}}
		
		\DecisionElementsTable
{Offer a basic membership plan; Offer a premium membership plan}
{High monthly demand; Low monthly demand}
{Monthly market conditions revealed after the membership-plan decision}
{Operating the selected activity under the realised conditions leaves funds available or a shortfall to cover. The monetary measure is profit in thousand USD for each plan--market-state pair.}

		
		\begin{CalculationLandscape}{14in}{11in}
\phantomsection
		\label{subsec:fitness-membership-c1-c3-c3-profit-table}
		\subsection*{C3 - Profit and payoff table \score{5}}
		
		\secInicio[sec:solutions-extraction][sec:fitness-membership-c1-c3]
		

\SymbolsTable{
$A_i$ & $A_1$: Offer a basic membership plan; $A_2$: Offer a premium membership plan & Alternative labels, in row order. \\
$S_j$ & $S_1$: High monthly demand; $S_2$: Low monthly demand & State labels, in column order; revealed after the choice. \\
$q(A_i,S_j)$ & Basic--high: 9; Basic--low: 6; Premium--high: 7; Premium--low: 4 hundreds of memberships & Membership lots for each plan under each demand level. \\
$r^A(A_i)$ & Basic plan: 8 thousand USD/100 memberships; Premium plan: 13 thousand USD/100 memberships & Plan-based revenue rate per hundred memberships. \\
$r^S(S_j)$ & High demand: 2 thousand USD/100 memberships; Low demand: 1 thousand USD/100 memberships & Additional revenue rate per hundred memberships from monthly demand. \\
$c^A(A_i)$ & Basic plan: 3 thousand USD/100 memberships; Premium plan: 6 thousand USD/100 memberships & Plan-based cost rate per hundred memberships. \\
$f_c(S_j)$ & High demand: 12 thousand USD; Low demand: 16 thousand USD & Fixed operating cost under each demand condition. \\
}

		
		\textbf{Profit equation.}
		\[
		R=q(A_i,S_j)\left[r^A(A_i)+r^S(S_j)\right],
		\qquad
		C=q(A_i,S_j)c^A(A_i)+f_c(S_j),
		\qquad
		P(A_i,S_j)=R-C.
		\]
		
		\CalculationPair{Revenue}{
\begin{tabular}{@{}lcc@{}}
\toprule Alternative & $S_1$ & $S_2$ \\ \midrule
$A_1$ & \SymbolicCell{q(A_1,S_1)\times \\ [r^A(A_1)+r^S(S_1)]}{R(A_1,S_1)} & \SymbolicCell{q(A_1,S_2)\times \\ [r^A(A_1)+r^S(S_2)]}{R(A_1,S_2)} \\[5pt]
$A_2$ & \SymbolicCell{q(A_2,S_1)\times \\ [r^A(A_2)+r^S(S_1)]}{R(A_2,S_1)} & \SymbolicCell{q(A_2,S_2)\times \\ [r^A(A_2)+r^S(S_2)]}{R(A_2,S_2)} \\[5pt]
\bottomrule
\end{tabular}
}{
\begin{tabular}{@{}lcc@{}}
\toprule Alternative & $S_1$ & $S_2$ \\ \midrule
$A_1$ & \CalcResult{9\times \\ [8+2]}{90} & \CalcResult{6\times \\ [8+1]}{54} \\[5pt]
$A_2$ & \CalcResult{7\times \\ [13+2]}{105} & \CalcResult{4\times \\ [13+1]}{56} \\[5pt]
\bottomrule
\end{tabular}
}
\CalculationPair{Cost}{
\begin{tabular}{@{}lcc@{}}
\toprule Alternative & $S_1$ & $S_2$ \\ \midrule
$A_1$ & \SymbolicCell{q(A_1,S_1)\times c^A(A_1) \\ {}+f_c(S_1)}{C(A_1,S_1)} & \SymbolicCell{q(A_1,S_2)\times c^A(A_1) \\ {}+f_c(S_2)}{C(A_1,S_2)} \\[5pt]
$A_2$ & \SymbolicCell{q(A_2,S_1)\times c^A(A_2) \\ {}+f_c(S_1)}{C(A_2,S_1)} & \SymbolicCell{q(A_2,S_2)\times c^A(A_2) \\ {}+f_c(S_2)}{C(A_2,S_2)} \\[5pt]
\bottomrule
\end{tabular}
}{
\begin{tabular}{@{}lcc@{}}
\toprule Alternative & $S_1$ & $S_2$ \\ \midrule
$A_1$ & \CalcResult{9\times 3 \\ {}+12}{39} & \CalcResult{6\times 3 \\ {}+16}{34} \\[5pt]
$A_2$ & \CalcResult{7\times 6 \\ {}+12}{54} & \CalcResult{4\times 6 \\ {}+16}{40} \\[5pt]
\bottomrule
\end{tabular}
}
\CalculationPair{Profit}{
\begin{tabular}{@{}lcc@{}}
\toprule Alternative & $S_1$ & $S_2$ \\ \midrule
$A_1$ & \SymbolicCell{R(A_1,S_1)-C(A_1,S_1)}{P(A_1,S_1)} & \SymbolicCell{R(A_1,S_2)-C(A_1,S_2)}{P(A_1,S_2)} \\[5pt]
$A_2$ & \SymbolicCell{R(A_2,S_1)-C(A_2,S_1)}{P(A_2,S_1)} & \SymbolicCell{R(A_2,S_2)-C(A_2,S_2)}{P(A_2,S_2)} \\[5pt]
\bottomrule
\end{tabular}
}{
\begin{tabular}{@{}lcc@{}}
\toprule Alternative & $S_1$ & $S_2$ \\ \midrule
$A_1$ & \CalcResult{90-39}{51} & \CalcResult{54-34}{20} \\[5pt]
$A_2$ & \CalcResult{105-54}{51} & \CalcResult{56-40}{16} \\[5pt]
\bottomrule
\end{tabular}
}

		
		\begingroup\renewcommand{\CalculationStateCount}{2}
\FinalPayoffTable{profit in thousand USD}{High monthly demand ($S_1$) & Low monthly demand ($S_2$)}{
Basic membership plan ($A_1$) & 51 & 20 \\
Premium membership plan ($A_2$) & 51 & 16 \\
}
\endgroup
\end{CalculationLandscape}

		
	}
	
	\sectionScored
	{Restaurant Expansion Strategy}
	{C4-\sml{2}}
	{sec:restaurant-expansion-c4}
	{
		A restaurant company must choose a strategy for expanding into a new
		location. Customer demand could be high, medium, or low. All payoffs are in
		thousands of dollars.
		
		\begin{center}
			\renewcommand{\arraystretch}{1.3}
			\small
			\begin{tabular}{lccc}
				\toprule
				& \multicolumn{3}{c}{Demand} \\
				\cmidrule(lr){2-4}
				Alternative & High & Medium & Low \\
				\midrule
				Full-service restaurant & 78 & 35 & 8 \\
				Casual restaurant & 72 & 48 & 32 \\
				Express restaurant & 95 & 42 & 28 \\
				Postpone expansion & 25 & 25 & 25 \\
				\bottomrule
			\end{tabular}
		\end{center}
		\vspace{-5mm}
		
		\textbf{Questions/Tasks.}
		\begin{enumerate}
			\CriterionTask
			{subsec:restaurant-expansion-c4}{4}
			{What would the optimistic decision be (Maximax)? What would the conservative decision be (Maximin)? \score{2}}
		\end{enumerate}
		
		\phantomsection
		\label{subsec:restaurant-expansion-c4}
		\subsection*{C4 - Maximax and Maximin \score{2}}
		
		\secInicio[sec:solutions-extraction][sec:restaurant-expansion-c4]
		
		\begin{multicols}{2}
			
			\textbf{Complete the table using letters first.}
			
			\[
			\begin{array}{c|ccccc}
				& A_1 & A_2 & A_3 & A_4 & \text{Decision}\\
				\hline
				\text{Max} &
				M_1 & M_2 & M_3 & M_4 &
				A^*_{\mathrm{Maximax}}=A_3
				\\
				\text{Min} &
				m_1 & m_2 & m_3 & m_4 &
				A^*_{\mathrm{Maximin}}=A_2
			\end{array}
			\]
			
			\textbf{Maximax}
			
			\[
			\begin{array}{ll}
				M_1 &= \max\{78,35,8\}=78\\
				M_2 &= \max\{72,48,32\}=72\\
				M_3 &= \max\{95,42,28\}=95\\
				M_4 &= \max\{25,25,25\}=25
			\end{array}
			\]
			
			\textbf{Maximin}
			
			\[
			\begin{array}{ll}
				m_1 &= \min\{78,35,8\}=8\\
				m_2 &= \min\{72,48,32\}=32\\
				m_3 &= \min\{95,42,28\}=28\\
				m_4 &= \min\{25,25,25\}=25
			\end{array}
			\]
			
			\columnbreak
			
			\textbf{Replace the letters with their values.}
			
			\[
			\begin{array}{c|ccccc}
				& A_1 & A_2 & A_3 & A_4 & \text{Decision}\\
				\hline
				\text{Max} &
				78 & 72 & 95 & 25 &
				A^*_{\mathrm{Maximax}}=A_3
				\\
				\text{Min} &
				8 & 32 & 28 & 25 &
				A^*_{\mathrm{Maximin}}=A_2
			\end{array}
			\]
			
			Therefore, the optimistic decision is
			\[
			A^*_{\mathrm{Maximax}}=A_3
			\quad\text{(Express restaurant)}
			\]
			
			and the conservative decision is
			\[
			A^*_{\mathrm{Maximin}}=A_2
			\quad\text{(Casual restaurant)}.
			\]
			
		\end{multicols}
		
	}
	
	\sectionScored
	{Hotel Development Strategy}
	{C4-\sml{2}}
	{sec:hotel-development-c4}
	{
		A hotel company must choose a strategy for developing a new property.
		Tourist demand could be high, medium, or low. All payoffs are in
		thousands of dollars.
		
		\begin{center}
			\renewcommand{\arraystretch}{1.3}
			\small
			\begin{tabular}{lccc}
				\toprule
				& \multicolumn{3}{c}{Tourist demand} \\
				\cmidrule(lr){2-4}
				Alternative & High & Medium & Low \\
				\midrule
				Large hotel & 72 & 34 & 10 \\
				Medium hotel & 65 & 46 & 22 \\
				Small hotel & 95 & 44 & 31 \\
				No development & 58 & 40 & 35 \\
				\bottomrule
			\end{tabular}
		\end{center}
		\vspace{-5mm}
		
		\textbf{Questions/Tasks.}
		\begin{enumerate}
			\CriterionTask
			{subsec:hotel-development-c4}{4}
			{What would the optimistic decision be (Maximax)? What would the conservative decision be (Maximin)? \score{2}}
		\end{enumerate}
		
		\phantomsection
		\label{subsec:hotel-development-c4}
		\subsection*{C4 - Maximax and Maximin \score{2}}
		
		\secInicio[sec:solutions-extraction][sec:hotel-development-c4]
		
		\begin{multicols}{2}
			
			\textbf{Complete the table using letters first.}
			
			\[
			\begin{array}{c|ccccc}
				& A_1 & A_2 & A_3 & A_4 & \text{Decision}\\
				\hline
				\text{Max} &
				M_1 & M_2 & M_3 & M_4 &
				A^*_{\mathrm{Maximax}}=A_3
				\\
				\text{Min} &
				m_1 & m_2 & m_3 & m_4 &
				A^*_{\mathrm{Maximin}}=A_4
			\end{array}
			\]
			
			\textbf{Maximax}
			
			\[
			\begin{array}{ll}
				M_1 &= \max\{72,34,10\}=72\\
				M_2 &= \max\{65,46,22\}=65\\
				M_3 &= \max\{95,44,31\}=95\\
				M_4 &= \max\{58,40,35\}=58
			\end{array}
			\]
			
			\textbf{Maximin}
			
			\[
			\begin{array}{ll}
				m_1 &= \min\{72,34,10\}=10\\
				m_2 &= \min\{65,46,22\}=22\\
				m_3 &= \min\{95,44,31\}=31\\
				m_4 &= \min\{58,40,35\}=35
			\end{array}
			\]
			
			\columnbreak
			
			\textbf{Replace the letters with their values.}
			
			\[
			\begin{array}{c|ccccc}
				& A_1 & A_2 & A_3 & A_4 & \text{Decision}\\
				\hline
				\text{Max} &
				72 & 65 & 95 & 58 &
				A^*_{\mathrm{Maximax}}=A_3
				\\
				\text{Min} &
				10 & 22 & 31 & 35 &
				A^*_{\mathrm{Maximin}}=A_4
			\end{array}
			\]
			
			Therefore, the optimistic decision is
			\[
			A^*_{\mathrm{Maximax}}=A_3
			\quad\text{(Small hotel)}
			\]
			
			and the conservative decision is
			\[
			A^*_{\mathrm{Maximin}}=A_4
			\quad\text{(No development)}.
			\]
			
		\end{multicols}
	}
		
	\sectionScored
	{Agricultural Technology Investment Strategy}
	{C4-\sml{2}}
	{sec:agricultural-technology-c4}
	{
		An agricultural company must choose a strategy for introducing a new
		automated farming system. Market conditions could be highly favourable,
		favourable, stable, or unfavourable. All payoffs are in thousands of dollars.
		
		\begin{center}
			\renewcommand{\arraystretch}{1.3}
			\small
			\begin{tabular}{lcccc}
				\toprule
				& \multicolumn{4}{c}{Market conditions} \\
				\cmidrule(lr){2-5}
				Alternative & Highly favourable & Favourable & Stable & Unfavourable \\
				\midrule
				Full automation & 125 & 58 & 8 & -42 \\
				Partial automation & 88 & 64 & 30 & 4 \\
				Basic automation & 62 & 52 & 38 & 20 \\
				Equipment leasing & 48 & 43 & 35 & 28 \\
				\bottomrule
			\end{tabular}
		\end{center}
		\vspace{-5mm}
		
		\textbf{Questions/Tasks.}
		\begin{enumerate}
			\CriterionTask
			{subsec:agricultural-technology-c4}{4}
			{What would the optimistic decision be (Maximax)? What would the conservative decision be (Maximin)? \score{2}}
		\end{enumerate}
		
		\phantomsection
		\label{subsec:agricultural-technology-c4}
		\subsection*{C4 - Maximax and Maximin \score{2}}
		
		\secInicio[sec:solutions-extraction][sec:agricultural-technology-c4]
		
		\begin{multicols}{2}
			
			\textbf{Complete the table using letters first.}
			
			\[
			\begin{array}{c|ccccc}
				& A_1 & A_2 & A_3 & A_4 & \text{Decision}\\
				\hline
				\text{Max} &
				M_1 & M_2 & M_3 & M_4 &
				A^*_{\mathrm{Maximax}}=A_1
				\\
				\text{Min} &
				m_1 & m_2 & m_3 & m_4 &
				A^*_{\mathrm{Maximin}}=A_4
			\end{array}
				\]
				
				\textbf{Maximax}
				
				\[
				\begin{array}{ll}
					M_1 &= \max\{125,58,8,-42\}=125\\
					M_2 &= \max\{88,64,30,4\}=88\\
					M_3 &= \max\{62,52,38,20\}=62\\
					M_4 &= \max\{48,43,35,28\}=48
				\end{array}
				\]
				
				\textbf{Maximin}
				
				\[
				\begin{array}{ll}
					m_1 &= \min\{125,58,8,-42\}=-42\\
					m_2 &= \min\{88,64,30,4\}=4\\
					m_3 &= \min\{62,52,38,20\}=20\\
					m_4 &= \min\{48,43,35,28\}=28
				\end{array}
				\]
				
				\columnbreak
				
				\textbf{Replace the letters with their values.}
				
				\[
				\begin{array}{c|ccccc}
					& A_1 & A_2 & A_3 & A_4 & \text{Decision}\\
					\hline
					\text{Max} &
					125 & 88 & 62 & 48 &
					A^*_{\mathrm{Maximax}}=A_1
					\\
					\text{Min} &
					-42 & 4 & 20 & 28 &
					A^*_{\mathrm{Maximin}}=A_4
				\end{array}
				\]
					
					Therefore, the optimistic decision is
					\[
					A^*_{\mathrm{Maximax}}=A_1
					\quad\text{(Full automation)}
					\]
					
					and the conservative decision is
					\[
					A^*_{\mathrm{Maximin}}=A_4
					\quad\text{(Equipment leasing)}.
					\]
					
				\end{multicols}
			}
			
	\sectionScored
	{Summarized Solutions}
	{C1-C4}
	{sec:summary}
	{
	\subsection{Local Food Market Stall \hfill C1--1 \quad C2--2}
		\begin{center}
			\begin{minipage}[t]{0.6\textwidth}
				\begin{tabularx}{\textwidth}
					{@{}p{1.5cm}>{\raggedright\arraybackslash}X@{}}
					\toprule
					Element & Description \\
					\midrule
					Alt. &
					\textbf{A1:} Standard stall \quad
					\textbf{A2:} Deluxe stall \\
					
					States &
					\textbf{S1:} High demand \quad
					\textbf{S2:} Low demand \\
					
					Events &
					Actual market demand after the stall decision \\
					
					Cons. &
					Profit in thousand US dollars for each alternative--state pair \\
					\bottomrule
				\end{tabularx}
			\end{minipage}
			\hfill
			\begin{minipage}[t]{0.35\textwidth}
				\begin{center}
					\begin{tabular}{lcc}
						\toprule
						Alt. &  \(\mathbf{S_1}\) & \(\mathbf{S_2}\)\\
						\midrule
						\( \mathbf{A_1}\) & 18 & 11\\
						\(\mathbf{A_2}\) & 31 & 6\\
						\bottomrule
					\end{tabular}
				\end{center}
			\end{minipage}
		\end{center}
	
	\subsection{Small Bakery Production \hfill C1--1 \quad C2--3}
	\begin{center}
		\begin{minipage}[t]{0.55\textwidth}
			\begin{tabularx}{\textwidth}{
					@{}p{1.5cm}>{\raggedright\arraybackslash}X@{}
				}
				\toprule
				Element & Description \\
				\midrule
				Alt. &
				\textbf{A1:} Standard oven \quad
				\textbf{A2:} Industrial oven \\
				
				States &
				\textbf{S1:} High demand \quad
				\textbf{S2:} Low demand \\
				
				Events &
				Actual customer demand after the oven decision \\
				
				Cons. &
				Profit in thousand BRL for each alternative--state pair \\
				\bottomrule
			\end{tabularx}
		\end{minipage}
		\hfill
		\begin{minipage}[t]{0.4\textwidth}
			\begin{tabular}{lcc}
				\toprule
				Alt. & \(\mathbf{S_1}\) & \(\mathbf{S_2}\)\\
				\midrule
				\( \mathbf{A_1}\)
				&
				\begin{tabular}{c}
					$16(3)+6$\\
					$=54$
				\end{tabular}
				&
				\begin{tabular}{c}
					$11(3)-3$\\
					$=30$
				\end{tabular}
				\\
				
				\( \mathbf{A_2}\)
				&
				\begin{tabular}{c}
					$13(5)+6$\\
					$=71$
				\end{tabular}
				&
				\begin{tabular}{c}
					$9(5)-3$\\
					$=42$
				\end{tabular}
				\\
				\bottomrule
			\end{tabular}
		\end{minipage}
	\end{center}
	
	\subsection{Urban Meal Subscription Service \hfill C1--2 \quad C2--5}
	\begin{center}
		\begin{minipage}[t]{0.55\textwidth}
			\begin{tabularx}{\textwidth}{
					@{}p{1.5cm}>{\raggedright\arraybackslash}X@{}
				}
				\toprule
				Element & Description \\
				\midrule
				Alt. &
				\textbf{A1:} Basic meal plan \quad
				\textbf{A2:} Premium meal plan \\
				
				States &
				\textbf{S1:} Strong customer interest \quad
				\textbf{S2:} Weak customer interest \\
				
				Events &
				Actual customer interest after the meal plan is selected \\
				
				Cons. &
				Profit in thousand BRL for each alternative--state pair \\
				\bottomrule
			\end{tabularx}
		\end{minipage}
		\hfill
		\begin{minipage}[t]{0.4\textwidth}
			\begin{tabular}{lcc}
				\toprule
				Alt. & \(\mathbf{S_1}\) & \(\mathbf{S_2}\)\\
				\midrule
				\( \mathbf{A_1}\)
				&
				\begin{tabular}{c}
					$13(4+2)+6$\\
					$=84$
				\end{tabular}
				&
				\begin{tabular}{c}
					$8(4+1)+6$\\
					$=46$
				\end{tabular}
				\\
				
				\( \mathbf{A_2}\)
				&
				\begin{tabular}{c}
					$10(6+2)+9$\\
					$=89$
				\end{tabular}
				&
				\begin{tabular}{c}
					$6(6+1)+9$\\
					$=51$
				\end{tabular}
				\\
				\bottomrule
			\end{tabular}
			
		\end{minipage}
	\end{center}
	
	\subsection{Regional Bookstore Expansion \hfill C1--1 \quad C3--2}
	\begin{center}
		\begin{minipage}[t]{0.55\textwidth}
			\begin{tabularx}{\textwidth}{
					@{}p{1.5cm}>{\raggedright\arraybackslash}X@{}
				}
				\toprule
				Element & Description \\
				\midrule
				Alt. &
				\textbf{A1:} Open children's section \quad
				\textbf{A2:} Maintain current selection \\
				
				States &
				\textbf{S1:} Strong traffic \quad
				\textbf{S2:} Weak traffic \\
				
				Events &
				Actual customer traffic after the product selection decision \\
				
				Cons. &
				Profit in million BRL for each alternative--state pair \\
				\bottomrule
			\end{tabularx}
		\end{minipage}
		\hfill
		\begin{minipage}[t]{0.4\textwidth}
			\begin{tabular}{lcc}
				\toprule
				Alt. & \(\mathbf{S_1}\) & \(\mathbf{S_2}\)\\
				\midrule
				\( \mathbf{A_1}\)
				&
				\begin{tabular}{c}
					$72-24$\\
					$=48$
				\end{tabular}
				&
				\begin{tabular}{c}
					$35-22$\\
					$=13$
				\end{tabular}
				\\
				
				\( \mathbf{A_2}\)
				&
				\begin{tabular}{c}
					$61-25$\\
					$=36$
				\end{tabular}
				&
				\begin{tabular}{c}
					$47-23$\\
					$=24$
				\end{tabular}
				\\
				\bottomrule
			\end{tabular}
			
		\end{minipage}
	\end{center}
	
	\subsection{Greenhouse Crop Production \hfill C1--1 \quad C3--3}
	\begin{center}
		\begin{minipage}[t]{0.55\textwidth}
			\begin{tabularx}{\textwidth}{
					@{}p{1.5cm}>{\raggedright\arraybackslash}X@{}
				}
				\toprule
				Element & Description \\
				\midrule
				Alt. &
				\textbf{A1:} Produce tomatoes \quad
				\textbf{A2:} Produce peppers \\
				
				States &
				\textbf{S1:} Favourable weather \quad
				\textbf{S2:} Unfavourable weather \\
				
				Events &
				Weather conditions revealed after the production decision \\
				
				Cons. &
				Profit in thousand EUR for each crop--weather-state pair \\
				\bottomrule
			\end{tabularx}
		\end{minipage}
		\hfill
		\begin{minipage}[t]{0.4\textwidth}
			\begin{tabular}{lcc}
				\toprule
				Alt. & \(\mathbf{S_1}\) & \(\mathbf{S_2}\)\\
				\midrule
				\( \mathbf{A_1}\)
				&
				\begin{tabular}{c}
					$8(18)-8(7)$\\
					$=88$
				\end{tabular}
				&
				\begin{tabular}{c}
					$5(18)-5(10)$\\
					$=40$
				\end{tabular}
				\\
				
				\( \mathbf{A_2}\)
				&
				\begin{tabular}{c}
					$7(15)-7(7)$\\
					$=56$
				\end{tabular}
				&
				\begin{tabular}{c}
					$4(15)-4(10)$\\
					$=20$
				\end{tabular}
				\\
				\bottomrule
			\end{tabular}
			
			
		\end{minipage}
	\end{center}
	
	\subsection{Fitness Center Membership Plans \hfill C1--2 \quad C3--5}
	{
	\begin{center}
		\begin{minipage}[t]{0.45\textwidth}
			\begin{tabularx}{\textwidth}{
					@{}p{1.5cm}>{\raggedright\arraybackslash}X@{}
				}
				\toprule
				Element & Description \\
				\midrule
				Alt. &
				\textbf{A1:} Offer a basic membership plan \quad
				\textbf{A2:} Offer a premium membership plan \\
				
				States &
				\textbf{S1:} High monthly demand \qquad \qquad \quad
				\textbf{S2:} Low monthly demand \\
				
				Events &
				Monthly market conditions revealed after the membership-plan decision \\
				
				Cons. &
				Profit in thousand USD for each plan--market-state pair \\
				\bottomrule
			\end{tabularx}
		\end{minipage}
		\hfill
		\begin{minipage}[t]{0.53\textwidth}
			\begin{tabular}{lcc}
				\toprule
				Alt. & \(\mathbf{S_1}\) & \(\mathbf{S_2}\)\\
				\midrule
				\( \mathbf{A_1}\)
				&
				\begin{tabular}{c}
					$9(8+2)-[9(3)+12]$\\
					$=51$
				\end{tabular}
				&
				\begin{tabular}{c}
					$6(8+1)-[6(3)+16]$\\
					$=20$
				\end{tabular}
				\\
				
				\( \mathbf{A_2}\)
				&
				\begin{tabular}{c}
					$7(13+2)-[7(6)+12]$\\
					$=51$
				\end{tabular}
				&
				\begin{tabular}{c}
					$4(13+1)-[4(6)+16]$\\
					$=16$
				\end{tabular}
				\\
				\bottomrule
			\end{tabular}
			
			
		\end{minipage}
	\end{center}
	}
	
	
	\clearpage
\begin{multicols}{3}
	\subsection{Restaurant Expansion Strategy  \hfill C4--2}
	\[
	\begin{array}{c|cc}
		& \text{Max} & \text{Min}\\
		\hline
		A_1 & 78 & 8\\
		A_2 & 72 & 32\\
		A_3 & 95 & 28\\
		A_4 & 25 & 25\\
		\hline
		A^*_{\mathrm{Maximax}} &
		A_3 & \\
		A^*_{\mathrm{Maximin}} & & A_2 
	\end{array}
	\]
	
	\subsection{Hotel Development Strategy \hfill C4--2}
	\[
	\begin{array}{c|cc}
		& \text{Max} & \text{Min}\\
		\hline
		A_1 & 72 & 10\\
		A_2 & 65 & 22\\
		A_3 & 95 & 31\\
		A_4 & 58 & 35\\
		\hline
		A^*_{\mathrm{Maximax}} &
		A_3 & \\
		A^*_{\mathrm{Maximin}} & & A_4
	\end{array}
	\]
	
	\subsection{Agricultural Technology Investment Strategy \hfill C4--2}
	\[
	\begin{array}{c|cc}
		& \text{Max} & \text{Min}\\
		\hline
		A_1 & 125 & -42\\
		A_2 & 88 & 4\\
		A_3 & 62 & 20\\
		A_4 & 48 & 28\\
		\hline
		A^*_{\mathrm{Maximax}} &
		A_1 & \\
		A^*_{\mathrm{Maximin}} & & A_4
	\end{array}
	\]
	
	\end{multicols}
	
\end{document}
